\documentclass[11pt,a4paper]{article}

\usepackage[utf8]{inputenc}
\usepackage[T1]{fontenc}
\usepackage[a4paper,margin=2cm]{geometry}
\renewcommand{\contentsname}{Sum\'ario}
\usepackage{lmodern}
\usepackage{microtype}
\usepackage[table]{xcolor}
\usepackage{colortbl}
\usepackage{array}
\usepackage{longtable}
\usepackage{ragged2e}
\usepackage{listings}
\usepackage{amssymb}
\usepackage[hidelinks,colorlinks=true,linkcolor=azulesc,urlcolor=azulesc]{hyperref}

% ---- Cores ----
\definecolor{azulesc}{HTML}{0B3D66}
\definecolor{azulmed}{HTML}{155080}
\definecolor{codebg}{HTML}{F4F6F8}
\definecolor{linhagrid}{HTML}{9BB7CC}
\definecolor{verdebox}{HTML}{2E7D32}
\definecolor{verdebg}{HTML}{EEF7EF}

% ---- Titulos de secao (sem titlesec) ----
\makeatletter
\renewcommand\section{\@startsection{section}{1}{0pt}%
  {-3.5ex \@plus -1ex \@minus -.2ex}{2.3ex \@plus.2ex}%
  {\normalfont\Large\bfseries\color{azulesc}}}
\renewcommand\subsection{\@startsection{subsection}{2}{0pt}%
  {-3.25ex\@plus -1ex \@minus -.2ex}{1.5ex \@plus .2ex}%
  {\normalfont\large\bfseries\color{azulmed}}}
\makeatother

% ---- Codigo ----
\lstset{
  basicstyle=\ttfamily\footnotesize,
  backgroundcolor=\color{codebg},
  frame=single,
  rulecolor=\color{linhagrid},
  breaklines=true,
  columns=fullflexible,
  keepspaces=true,
  showstringspaces=false,
  xleftmargin=4pt,
  framexleftmargin=4pt,
  aboveskip=8pt,
  belowskip=8pt
}

% ---- Colunas de tabela ----
\newcolumntype{L}[1]{>{\RaggedRight\arraybackslash}p{#1}}
\newcommand{\hc}[1]{\textcolor{white}{\textbf{#1}}}
\arrayrulecolor{linhagrid}
\renewcommand{\arraystretch}{1.25}

% ---- Caixa de destaque ----
\newcommand{\nota}[1]{%
  \begin{center}\colorbox{verdebg}{\parbox{0.95\linewidth}{%
  \color{verdebox}#1}}\end{center}}

\setlength{\parindent}{0pt}
\setlength{\parskip}{5pt}

\begin{document}

% ===================== CAPA =====================
\begin{titlepage}
\vspace*{3cm}
\begin{center}
{\color{azulesc}\Huge\bfseries SMART-JBMC}\\[10pt]
{\color{azulmed}\LARGE Guia de Instala\c{c}\~ao e Configura\c{c}\~ao}\\[40pt]
\rule{0.7\linewidth}{2pt}\\[20pt]
{\large Plataforma de verifica\c{c}\~ao formal (JBMC + ESBMC)}\\[6pt]
{\large com classifica\c{c}\~ao CWE/CERT-C e reparo guiado por contraexemplo}\\[40pt]
\begin{tabular}{rl}
\textbf{Disciplina:} & Verifica\c{c}\~ao de Software (Mestrado)\\[4pt]
\textbf{Orientador:} & Prof. Lucas Cordeiro (University of Manchester)\\[4pt]
\textbf{Vers\~ao:} & 0.1.0\\[4pt]
\textbf{Data:} & 23 de maio de 2026\\
\end{tabular}
\end{center}
\vfill
\end{titlepage}

\tableofcontents
\newpage

% ===================== 1 =====================
\section{Vis\~ao geral}
O SMART-JBMC \'e uma plataforma web (backend FastAPI + front-end SPA) que orquestra
verificadores formais para Java, Kotlin e C, classifica viola\c{c}\~oes em CWE/CERT-C
e (em etapa futura) repara bugs com aux\'ilio de LLM. Este guia cobre a instala\c{c}\~ao
das ferramentas externas, do projeto e da configura\c{c}\~ao, al\'em de como levar o
projeto para outra m\'aquina.

\nota{\textbf{Resolu\c{c}\~ao de ferramentas em 3 n\'iveis.} Cada ferramenta (esbmc, jbmc,
kotlinc, \dots) \'e localizada nesta ordem de prioridade: \textbf{(1)} caminho no
\texttt{.env} \(\to\) \textbf{(2)} bin\'ario copiado em \texttt{tools/} \(\to\)
\textbf{(3)} \texttt{PATH} do sistema. Isso torna o projeto port\'avel entre m\'aquinas.}

% ===================== 2 =====================
\section{Requisitos de sistema}
Sistema operacional Linux x86\_64 (testado em Ubuntu 24.04). Python 3.10 ou superior.

\begin{center}
\footnotesize
\begin{longtable}{|L{2.6cm}|L{5.4cm}|L{3cm}|L{2.2cm}|}
\hline
\rowcolor{azulesc} \hc{Ferramenta} & \hc{Para qu\^e} & \hc{Vers\~ao testada} & \hc{Obrigat\'oria?} \\ \hline
\endhead
\texttt{python3} & Runtime do backend & 3.12 & Sim \\ \hline
\texttt{esbmc} & Verifica\c{c}\~ao de C & 8.2.0 & Sim (alvos C) \\ \hline
\texttt{jbmc} & Verifica\c{c}\~ao de Java/Kotlin (suite CBMC) & 6.9.0 & Sim (alvos JVM) \\ \hline
\texttt{javac} & Compilar Java para bytecode & JDK 21 & Sim (alvos Java) \\ \hline
\texttt{java} & Runtime JVM / rodar testes & JDK 21 & Sim (alvos JVM) \\ \hline
\texttt{kotlinc} & Compilar Kotlin para bytecode & 2.3.21 & S\'o p/ Kotlin \\ \hline
\texttt{gcc} & Compilar testes C & qualquer recente & S\'o p/ testes C \\ \hline
\end{longtable}
\end{center}

% ===================== 3 =====================
\section{Instalando as ferramentas externas}

\subsection{ESBMC (verifica\c{c}\~ao de C)}
Baixe o bin\'ario oficial em \url{https://github.com/esbmc/esbmc/releases} (release
Linux). Descompacte e anote o caminho do execut\'avel, por exemplo
\texttt{/home/usuario/esbmc/bin/esbmc}. Verifique:
\begin{lstlisting}[language=bash]
/home/usuario/esbmc/bin/esbmc --version
# ESBMC version 8.2.0 64-bit x86_64 linux
\end{lstlisting}

\subsection{JBMC (verifica\c{c}\~ao de Java/Kotlin)}
O JBMC faz parte da suite \textbf{CBMC} (Diffblue), distribu\'ida \emph{separadamente}
do ESBMC. Baixe a release em \url{https://github.com/diffblue/cbmc/releases}:
\begin{lstlisting}[language=bash]
# Ubuntu/Debian: instale o pacote .deb da release (inclui jbmc e cbmc)
sudo apt install ./ubuntu-24.04-cbmc-6.9.0-Linux.deb

# Verifique
jbmc --version       # 6.9.0 (cbmc-6.9.0)
\end{lstlisting}
\nota{O \texttt{jbmc} precisa de um JDK instalado (ver \texttt{javac}/\texttt{java}) para
compilar e verificar o bytecode.}

\subsection{Kotlin (kotlinc) via SDKMAN}
O \texttt{kotlinc} \'e um script que depende do SDK Kotlin e de um JRE. A forma mais
simples \'e via SDKMAN:
\begin{lstlisting}[language=bash]
curl -s "https://get.sdkman.io" | bash
source "$HOME/.sdkman/bin/sdkman-init.sh"
sdk install kotlin
kotlinc -version     # kotlinc-jvm 2.3.21 (JRE 21...)
\end{lstlisting}

\subsection{JDK (javac e java)}
\begin{lstlisting}[language=bash]
sudo apt install openjdk-21-jdk
javac -version && java -version
\end{lstlisting}

\subsection{GCC}
\begin{lstlisting}[language=bash]
sudo apt install build-essential   # fornece o gcc
gcc --version
\end{lstlisting}

% ===================== 4 =====================
\section{Instalando o projeto}
\begin{lstlisting}[language=bash]
cd ~/tools/TrbLuccas/smart-jbmc

# Ambiente virtual isolado (recomendado)
python3 -m venv .venv
source .venv/bin/activate

# Dependencias do backend
pip install -r requirements.txt
\end{lstlisting}

% ===================== 5 =====================
\section{Configura\c{c}\~ao (\texttt{.env})}
Copie o modelo e ajuste os caminhos para a sua m\'aquina:
\begin{lstlisting}[language=bash]
cp .env.example .env
\end{lstlisting}

Exemplo de \texttt{.env} (os caminhos s\~ao espec\'ificos de cada m\'aquina):
\begin{lstlisting}
ESBMC_PATH=/home/franklin/esbmc/bin/esbmc
JBMC_PATH=/usr/bin/jbmc
KOTLINC_PATH=/home/franklin/.sdkman/candidates/kotlin/current/bin/kotlinc

# Reparo guiado por LLM: so o modelo padrao. A chave e de cada usuario
DEEPSEEK_MODEL=deepseek-v4-flash
\end{lstlisting}

\begin{center}
\footnotesize
\begin{longtable}{|L{3.4cm}|L{10.6cm}|}
\hline
\rowcolor{azulesc} \hc{Vari\'avel} & \hc{Descri\c{c}\~ao} \\ \hline
\endhead
\texttt{ESBMC\_PATH} & Caminho do execut\'avel ESBMC. Vazio \(\to\) usa \texttt{tools/esbmc} ou o \texttt{PATH}. \\ \hline
\texttt{JBMC\_PATH} & Caminho do execut\'avel JBMC. \\ \hline
\texttt{KOTLINC\_PATH} & Caminho do \texttt{kotlinc} (s\'o para alvos Kotlin). \\ \hline
\texttt{JAVAC\_PATH}, \texttt{JAVA\_PATH}, \texttt{GCC\_PATH} & Opcionais; normalmente resolvidos pelo \texttt{PATH}. \\ \hline
\texttt{DEEPSEEK\_API\_KEY} & Chave da API DeepSeek (s\'o para o reparo-LLM). Nunca versione. \\ \hline
\texttt{DEEPSEEK\_MODEL} & \texttt{deepseek-v4-flash} (econ\^omico) ou \texttt{deepseek-v4-pro}. \\ \hline
\end{longtable}
\end{center}

\nota{\textbf{Seguran\c{c}a:} o arquivo \texttt{.env} cont\'em caminhos locais e a chave de
API. Ele j\'a est\'a no \texttt{.gitignore} --- \textbf{nunca} o versione nem o
compartilhe.}

% ===================== 6 =====================
\section{Executando}
\begin{lstlisting}[language=bash]
cd ~/tools/TrbLuccas/smart-jbmc
.venv/bin/uvicorn smart_jbmc.api.app:app --app-dir backend --host 127.0.0.1 --port 8000
\end{lstlisting}
Abra no navegador:
\begin{itemize}
  \item \url{http://localhost:8000} --- a SPA (painel, estado do servi\c{c}o)
  \item \url{http://localhost:8000/api/docs} --- documenta\c{c}\~ao autom\'atica da API (Swagger)
\end{itemize}

% ===================== 7 =====================
\section{Verificando a instala\c{c}\~ao}
O endpoint \texttt{/api/health} reporta a disponibilidade e o caminho de cada
ferramenta:
\begin{lstlisting}[language=bash]
curl -s http://localhost:8000/api/health
\end{lstlisting}
Cada ferramenta deve aparecer com \texttt{"available": true}. No painel web, os
\emph{badges} ficam verdes e o caminho aparece ao passar o mouse. Para rodar os testes:
\begin{lstlisting}[language=bash]
.venv/bin/python -m pytest -q
\end{lstlisting}

% ===================== USO =====================
\section{Uso da interface e par\^ametros}
Abra \url{http://localhost:8000}. O card \textbf{Ambiente} (no topo) indica quais
ferramentas foram encontradas. Em \textbf{Nova verifica\c{c}\~ao}:
\begin{itemize}
  \item \textbf{Pasta} --- digite o caminho ou clique em \textbf{Procurar} para navegar
        pelos diret\'orios da m\'aquina (mostra subpastas e quantos arquivos-fonte h\'a).
  \item \textbf{Verificar} --- uma barra de progresso mostra a fase (escanear $\to$
        compilar $\to$ verificar) e os alvos conclu\'idos em tempo real.
  \item \textbf{Resultados} --- resumo (BUG/SAFE/TIMEOUT/ERROR + CASTLE Score),
        distribui\c{c}\~ao CWE e tabela de alvos; clique num alvo para ver o contraexemplo.
  \item \textbf{Hist\'orico} --- lista as runs; o item ativo fica destacado e o cabe\c{c}alho
        identifica qual run est\'a aberta.
\end{itemize}

\begin{center}
\footnotesize
\begin{longtable}{|L{2.3cm}|L{1.3cm}|L{1.9cm}|L{8cm}|}
\hline
\rowcolor{azulesc} \hc{Par\^ametro} & \hc{M\'in} & \hc{M\'ax} & \hc{O que faz / implica\c{c}\~ao} \\ \hline
\endhead
\texttt{unwind} & 1 & $\sim$100 & Bound do BMC (desenrolamento de loops). $\uparrow$ mais confian\c{c}a, por\'em mais lento e mais mem\'oria; $\downarrow$ r\'apido, mas pode perder bugs profundos (BMC \'e incompleto al\'em do bound). T\'ipico: 5--20. \\ \hline
\texttt{timeout} & 1\,s & 3600\,s & Tempo m\'aximo do solver por alvo. $\uparrow$ menos \texttt{TIMEOUT} em alvos dif\'iceis (run mais lenta); $\downarrow$ corta alvos pesados cedo. T\'ipico: 30--120. \\ \hline
\texttt{jobs} & 1 & n\'ucleos & Verifica\c{c}\~oes em paralelo. $\uparrow$ mais r\'apido em multi-core (usa mais CPU/RAM); \texttt{1} = sequencial (bom para depurar). \\ \hline
\texttt{solver} & --- & --- & Backend SMT do ESBMC: \texttt{boolector} (r\'apido, padr\~ao), \texttt{z3} (aritm\'etica), \texttt{mathsat} (interpola\c{c}\~ao). O JBMC usa MiniSat (independe desta op\c{c}\~ao). \\ \hline
\end{longtable}
\end{center}

% ===================== REPARO =====================
\section{Reparo guiado por contraexemplo (DeepSeek V4)}
Ap\'os uma verifica\c{c}\~ao com bugs, a interface mostra \textbf{Reparar bugs}. O ciclo:
a LLM (DeepSeek V4) prop\~oe um patch a partir do contraexemplo; o verificador
(JBMC/ESBMC) re-verifica; se o bug persiste, o \textbf{novo contraexemplo} realimenta
a LLM, at\'e o BMC provar \texttt{FIX\_OK} ou esgotar as itera\c{c}\~oes. O verificador
\'e o or\'aculo: s\'o conta como reparado o que ele prova seguro.

\begin{itemize}
  \item \textbf{Chave:} preencha \texttt{DEEPSEEK\_API\_KEY} no \texttt{.env} (e mantenha
        saldo na conta). \texttt{DEEPSEEK\_MODEL}: \texttt{deepseek-v4-flash} (econ\^omico,
        padr\~ao) ou \texttt{deepseek-v4-pro}.
  \item \textbf{Modo \texttt{dry-run}} (interruptor na tela): executa todo o loop
        \textbf{sem chamar a API} (n\~ao consome saldo) --- \'util para validar a fia\c{c}\~ao.
  \item \textbf{M\'etricas:} taxa de reparo verificado e itera\c{c}\~oes m\'edias por
        categoria, para comparar com o baseline heur\'istico.
  \item \textbf{REST:} \texttt{POST /api/repair} e \texttt{GET /api/repair/\{run\_id\}}.
\end{itemize}

% ===================== 8 =====================
\section{Portabilidade --- levar para outro computador}
O projeto foi desenhado para ser port\'avel. Recomenda\c{c}\~ao:
\begin{enumerate}
  \item Copie a pasta \texttt{smart-jbmc/} \textbf{sem} \texttt{.venv/}, \texttt{tools/} e \texttt{.env} (s\~ao espec\'ificos da m\'aquina).
  \item Na nova m\'aquina, instale as ferramentas externas (Se\c{c}\~ao 3) e o JDK.
  \item Recrie o \texttt{.env} com os caminhos locais (use o \texttt{.env.example}).
  \item \texttt{pip install -r requirements.txt} e execute (Se\c{c}\~ao 6).
\end{enumerate}

\nota{Se a outra m\'aquina for \textbf{Linux x86\_64 igual}, voc\^e pode copiar os
bin\'arios \texttt{esbmc} e \texttt{jbmc} para a pasta \texttt{tools/} do projeto --- eles
ser\~ao detectados automaticamente, sem precisar do \texttt{.env}. O \texttt{kotlinc},
por\'em, sempre exige o SDK Kotlin + JRE instalados na m\'aquina destino.}

% ===================== 9 =====================
\section{Solu\c{c}\~ao de problemas}
\begin{center}
\footnotesize
\begin{longtable}{|L{5cm}|L{9cm}|}
\hline
\rowcolor{azulesc} \hc{Sintoma} & \hc{Solu\c{c}\~ao} \\ \hline
\endhead
\texttt{jbmc} aparece como indispon\'ivel & Instale a suite CBMC (Se\c{c}\~ao 3.2) e/ou defina \texttt{JBMC\_PATH} no \texttt{.env}. \\ \hline
\texttt{kotlinc} indispon\'ivel & Instale via SDKMAN; aponte \texttt{KOTLINC\_PATH} para \texttt{.../kotlin/current/bin/kotlinc}. \\ \hline
``Address already in use'' na porta 8000 & Outro servidor j\'a est\'a rodando: \texttt{pkill -f "uvicorn smart\_jbmc"} ou use \texttt{--port 8001}. \\ \hline
SPA carrega mas mostra ``offline'' & O backend n\~ao subiu; confira o terminal do \texttt{uvicorn}. \\ \hline
Pylance n\~ao acha \texttt{fastapi} & Selecione o interpretador do \texttt{.venv} no editor. \\ \hline
\end{longtable}
\end{center}

\end{document}
